\section{From the human proofs to first-order \ipra}\label{sec:internalization}

This section is the proof-theoretic bridge.  Its purpose is to show that the preceding arguments are not merely constructive in an informal philosophical sense: after fixing their input names, every proof reduces to bounded arithmetic and quantifier-free induction.

\subsection{Finite-contract normal form}

Every analytic assertion used in the paper can be compiled to the following shape.  Let $u$ code all ordinary parameters, $p$ the requested precision, and $c$ a finite certificate.  There are primitive-recursive predicates
\[
 \operatorname{InputOK}(u,p),
 \qquad
 \operatorname{CertOK}(u,p,c),
 \qquad
 \operatorname{Conclusion}(u,p,c)
\]
such that the object theorem is
\begin{equation}\label{eq:contract-normal-form}
 \forall u\,\forall p\,\forall c\;
 \bigl(\operatorname{InputOK}(u,p)=1\wedge
       \operatorname{CertOK}(u,p,c)=1\bigr)
 \to \operatorname{Conclusion}(u,p,c)=1.
\end{equation}
When the certificate is produced by a primitive-recursive function $C(u,p)$, substitute it into \eqref{eq:contract-normal-form} and prove
$\operatorname{CertOK}(u,p,C(u,p))=1$.  Thus no existential choice is required at run time.

For example, a CDF enclosure certificate contains:
\begin{enumerate}[label=(\roman*)]
  \item a dyadic interval for $\pi$ and its Machin-series remainder checks;
  \item a dyadic square-root interval for $2\pi$;
  \item the finite rational sum in \eqref{eq:cdf-taylor};
  \item the next-term bound and a Boolean check that the terms are decreasing;
  \item outward-rounded lower and upper endpoints.
\end{enumerate}
Every field is a natural-number code.  The checker uses only bounded loops and cross-multiplied rational comparisons.

\subsection{Coding lists, matrices, polynomials, and intervals}

The existing ProveIt arithmetic library already supports pairing, bounded sums, division with remainder, gcd, and simultaneous recursion.  Lists can therefore be coded by a length and a mixed-radix or nested-pair payload.  The relevant operations are primitive recursive:
\begin{itemize}
  \item lookup with a default value;
  \item map, bounded fold, convolution, and coefficientwise equality;
  \item finite permutations and pairing enumerators;
  \item matrix lookup, transpose, product, determinant, and quadratic forms;
  \item interval addition, multiplication, reciprocal under apartness, and dyadic hull.
\end{itemize}
The object proofs of their defining equations use induction on the list length or loop counter.  Bounds needed for a mixed-radix code are explicit fields of the input contract, so no unbounded minimization is hidden.

\subsection{Names as finite tables}

A fast name through precision horizon $h$ is a table
\[
 A=(a_0,\dots,a_h)
\]
with the bounded regularity predicate
\begin{equation}\label{eq:finite-fast}
 \forall i,j\le h\quad
 |a_i-a_j|\le2^{-i}+2^{-j}.
\end{equation}
The universal statement about an infinite name is the externally indexed family of finite theorems for all $h$.  The table may be generated by a fixed \code{PR} term, in which case its entries do not have to be supplied in a certificate.  Arithmetic on names is implemented by backward precision allocation.  Multiplication requests magnitude bounds; reciprocal requests a lower apartness bound; composition requests the outer modulus before deciding the inner precision.

An equality of represented values at precision $p$ becomes a rational inequality between entries at an explicitly larger index.  A derivative contract is a bounded formula of the form
\begin{equation}\label{eq:derivative-contract}
 |F(x+h,r)-F(x,r)-hD(x,r)|
 \le |h|2^{-p}+2^{-r}C
\end{equation}
for rational grid points and bounded $h$.  An integral contract compares a finite Riemann sum with an output interval.  These are quantifier-free once all displayed bounded quantifiers are compiled to products of Boolean characteristic terms.

\subsection{Object-language proof of the numerical primitives}

The following lemmas are the minimal reusable analytic basis.

\begin{enumerate}[label=\textbf{N\arabic*.}]
  \item \textbf{Dyadic bisection.}  The integer loop preserves
  $L^2\le q2^{2p}<U^2$ and $U-L$ halves at every stage.  Quantifier-free induction proves the final square-root enclosure.
  \item \textbf{Alternating finite sums.}  If $0\le a_{j+1}\le a_j$, even and odd partial sums bracket every later partial sum.  The proof is induction on the number of paired terms.  Applying it to arctangent and $e^{-y}$ gives \eqref{eq:atan-bound} and \eqref{eq:exp-alt}.
  \item \textbf{Finite Taylor identity.}  Repeated multiplication proves the coefficient recursion and a telescoping formula for the remainder.  The Taylor packet already supplies a general contract; the Gaussian development needs only exponential and trigonometric specializations.
  \item \textbf{Finite Riemann refinement.}  Split every coarse cell into a fixed number of fine cells and compare sample values by the supplied modulus.  Bounded-sum triangle inequalities prove the mesh error.
  \item \textbf{Finite integration by parts.}  On a partition, summation by parts is an exact identity.  Derivative contracts bound the discrepancy from the integral formula.  The tail producer then removes the finite endpoints.
\end{enumerate}

Each lemma has a direct first-order formula because rational arithmetic and every loop bound are explicit.

\subsection{Moment and Hermite derivations}

The moment theorem is particularly close to a completed object proof.  Define \code{mom : PR 1} by simultaneous primitive recursion:
\[
 m(0)=1,
 \quad m(1)=0,
 \quad m(n+2)=(n+1)m(n).
\]
The target formula states that the represented integral moment is within $2^{-p}$ of the rational integer $m(n)$.  Its induction invariant contains the finite integration-by-parts identity at cutoff $T(n,p)$ and the two polynomial tail bounds.  The successor-by-two step is a finite term rewrite plus one use of the integration-by-parts lemma.  Parity and double factorial are already primitive recursive, so a second induction proves
$m(2r)=(2r-1)!!$ and $m(2r+1)=0$.

For Hermite polynomials, encode a coefficient list $H_n$.  The recurrence
\[
 H_{n+1}=XH_n-nH_{n-1}
\]
is a primitive-recursive list operation.  A simultaneous induction proves:
\begin{align*}
 \operatorname{deriv}(H_n)&=nH_{n-1},\\
 P_n&=(-1)^nH_n,
\end{align*}
where $P_0=1$ and $P_{n+1}=P_n'-XP_n$ is the polynomial multiplying the $n$th Gaussian derivative.  Equality is coefficientwise integer equality, so this entire theorem lies in the exact algebraic layer before any represented analysis is invoked.

\subsection{Transforms and independence}

A complex number is a pair of rational names.  Complex addition, multiplication, and exponential on bounded boxes are primitive-recursive.  A characteristic-function sample at rational $t$ is a pair of finite integrals of cosine and sine.  The proof of the Gaussian transform can be compiled in either of two ways:
\begin{enumerate}[label=(\alph*)]
  \item reuse the zeta packet's Gaussian transform contract and specialize its parameters; or
  \item prove the finite ODE residual for $C_T(t)=\int_{-T}^T\phi(x)\cos(tx)\dd x$, bound the two endpoint terms, and invoke a represented ODE uniqueness lemma derived from the finite fundamental theorem.
\end{enumerate}
The second route has fewer dependencies.  Multiplication of characteristic functions under independence is initially the exact equality of two finite product-space sums.  For density presentations it is the finite rectangle Fubini identity plus tail errors.  Linear-combination and multivariate theorems then reduce to polynomial identities in the exponent.

Independence itself is never treated as an unbounded set-theoretic predicate.  For finite rational laws it is a table factorization.  For represented densities, the theorem's conclusion is the family of rectangle-probability equalities to precision $2^{-p}$.  Transform factorization is converted to those finite observations by the same smoothing contract used in the central-limit packet.

\subsection{Information integrals}

The entropy and divergence theorems compile to a finite box $[-T,T]$ plus three certificates:
\begin{enumerate}[label=(\roman*)]
  \item a lower bound for every denominator or logarithm argument on the box;
  \item a quadrature and evaluation error budget on the box;
  \item a tail bound for the weighted logarithmic integrand.
\end{enumerate}
For Gaussian-versus-Gaussian formulas, completion of squares reduces the integral exactly to normalization and the first two moments.  For maximum entropy, the comparison density supplies the third certificate as a hypothesis.  Gibbs' inequality is compiled from the rational polynomial/logarithm enclosure of
$u\log u-u+1$ on dyadic subintervals; near $u=0$ the explicit bound $|u\log u|\le C\sqrt u$ handles the endpoint.  The nonnegativity conclusion is a rational lower enclosure, so intuitionistic logic is sufficient.

\subsection{Matrix and sample proofs}

For fixed dimension $d$, determinants, adjugates, and quadratic forms are bounded sums.  For variable finite $d$, list-coded Gaussian elimination carries a pivot-apartness certificate.  The multivariate theorem can avoid choosing eigenvectors: its primitive input is a factor $A$ together with a checked identity $AA^T=\Sigma$.  Orthogonal invariance is the finite identity $QQ^T=I$ inserted into the characteristic exponent.

The sample decomposition needs a matrix whose first row is the normalized all-ones vector.  There are two arithmetizations.
\begin{itemize}
  \item For each fixed $n$, compile an algebraic Householder matrix over the field generated by $\sqrt n$ and verify $QQ^T=I$ by finite algebra.
  \item Uniformly in $n$, avoid the full matrix and prove that the scalar mean coordinate and a coded orthogonal basis of the rational hyperplane have the required covariance, with square-root approximations carried relationally.
\end{itemize}
The first is shorter for a theorem family; the second is better for a single uniform library.  In both cases the crucial Pythagorean identity \eqref{eq:ss-decomp} is an exact ring theorem already within reach of ProveIt's derivation tools.

\subsection{Intuitionistic validity}

No proof uses excluded middle for an unbounded formula.  Apparent case distinctions are of four harmless kinds:
\begin{enumerate}[label=(\roman*)]
  \item bounded equality or order of integers/rationals, hence decidable by a primitive-recursive Boolean;
  \item a supplied apartness witness, such as $\sigma\ge2^{-s_0}$;
  \item an interval algorithm that returns a relation rather than deciding exact real equality;
  \item separate theorem schemas under separately stated hypotheses.
\end{enumerate}
The quantile algorithm is the main example: when comparison with $p$ is undecided at a midpoint, it returns a small interval using the positive derivative lower bound instead of invoking trichotomy for reals.

\subsection{Compilation theorem}

\begin{theorem}[iPRA compilation criterion]\label{thm:compilation}
Suppose a theorem in this manuscript has:
\begin{enumerate}[label=(\roman*)]
  \item input and output represented by fixed primitive-recursive tables;
  \item every division guarded by a displayed rational apartness witness;
  \item every limit, derivative, integral, or series guarded by an explicit primitive-recursive modulus or tail bound;
  \item a proof consisting of finite algebra, bounded sums/products/search, and inductions whose invariants are quantifier-free after bounded quantifiers are Booleanized.
\end{enumerate}
Then, for each fixed choice of the represented programs, its precision-indexed conclusion is derivable in ordinary \ipra.  After encoding the programs as \code{PR} terms, the derivation can be represented by a Lean term of type \code{[] $\vdash$[PRAAx] $\varphi$} and transported to the native concise calculus by \code{pra\_iff\_native}.
\end{theorem}

\begin{proof}
Primitive-recursive tables and checkers are available as function symbols with defining equations in \ipra.  Bounded quantifiers are finite Boolean folds.  Every finite algebraic step is a derivation from equality and order lemmas; every loop invariant is proved by quantifier-free induction.  Moduli replace all epsilon-limit quantifier alternations by explicit numerical substitutions.  Therefore the generated proof uses only \code{PRAAx}.  ProveIt's bi-interpretation theorem supplies the final transport.  The proof is a compilation recipe rather than an appeal to a semantic completeness theorem.
\end{proof}

Every result marked P or E in the ledger was deliberately stated to meet \cref{thm:compilation}.  The remaining work is engineering the explicit Lean proof terms, not discovering an additional classical principle.
